\addtotoc{Товчлол}
\thispagestyle{plain}

\vspace*{-6mm}
\begin{center}
\textbf{\large \titlename}\\[1mm]
\end{center}

\begin{flushright}
\departmentname\\
Оюутан: \studentID, \authorShort\\
Удирдагч багш: \supervisorname
\end{flushright}

\vspace{1mm}
\setcounter{table}{0}
\setcounter{figure}{0}
\renewcommand{\thetable}{\arabic{table}}
\renewcommand{\thefigure}{\arabic{figure}}
\setlength{\columnsep}{7mm}
\begin{multicols}{2}
\setlength{\parskip}{5pt}
\sloppy

\noindent\textbf{1. Удиртгал}

Энэхүү судалгаа нь валютын зах зээл (Forex)-д машин сургалтын загвараар дохио үүсгэж, серверийн хэсэг (backend) болон харагдах хэсэг (frontend)-тэй нэгдсэн систем хэлбэрээр хэрэгжүүлсэн. Санхүүгийн өгөгдөл нь шуугиан ихтэй, горимын өөрчлөлтөд мэдрэмтгий учраас өндөр нарийвчлалыг дангаар хөөхөөс илүү арга зүй-гүйцэтгэл-эрсдэлийн гурван түвшний баталгаажуулалтыг хэрэгжүүлсэн.

Судалгааны зорилт нь: (1) EUR/USD дээр тогтвортой ансамбль загвар бүтээх, (2) дохионы чанарыг бодит нөхцөлд үнэлэх, (3) прототип системийг хэрэглэгчид хүрэх түвшинд холбох, (4) өгөөж ба эрсдэлийг статистик аргаар тайлбарлах.

\noindent\textbf{2. Онолын суурь}

GBDT бүлгийн LightGBM, XGBoost, CatBoost алгоритмуудыг зөөлөн саналын (soft voting) ансамблиар нэгтгэхэд тулгуурласан. Хүснэгтэн өгөгдөлд гүн сургалтаас илүү тогтвортой ажиллах давуу талтай. Урагш гүйлгэн баталгаажуулалт (WFV) аргаар цаг хугацааны дарааллыг хадгалж, ирээдүйн мэдээлэл ашиглах хэвийдэл (look-ahead bias)-ийн эрсдэлийг бууруулсан.

\noindent\textbf{3. Арга зүй}

Дизайн шинжлэх ухааны ``бүтээх-үнэлэх-сайжруулах'' мөчлөгөөр 7 давталтаар хэрэгжүүлсэн. EUR/USD-ийн 2015--2024 оны OHLCV өгөгдлийг M1--H4 зургаан интервалаар ашиглаж, интервал бүрээс 8 шинж чанар (feature) гаргаж нийт 48 шинж чанарын матриц байгуулсан. Шошго (BUY/SELL/HOLD) нь ирээдүйн 240 минутын ханшийн хөдөлгөөнд суурилсан. Дохионы шүүлтүүрт итгэлцлийн босго ($\geq 0.90$) болон ATR ($\geq 4.0$ пипс)-ыг ашигласан. WFV-ийн Fold 5 огтлолд (2015--2022 / 2023 / 2024) KPI тайлан төвлөрүүлсэн.

Түүхэн өгөгдөл дээрх системийн туршилтыг 2025.01.01--2026.01.01 цонхонд, \$10,000 анхны хөрөнгөтэй, бодит спрэд нөхцөлд хийсэн. EA-ийн босго 0.90 дээр тогтмол байсан бөгөөд 0.90--0.99 мужийн ялгаа нь урьдчилан шүүсэн дохионы багцуудын сценарийн харьцуулалтаас үүдсэн.

\noindent\textbf{4. Үр дүн ба дүгнэлт}

359{,}639 таамгаас шүүлтүүрийн дараа 3{,}567 дохио (1.0\%) үлдсэн. Бүх босгын сценари эерэг өгөөжтэй гарсан: хамгийн өндөр нь 0.91 босгон дээр +517.70\% (\$61{,}770.46), 0.90 дээр +482.62\% (\$58{,}262.39). EUR/USD buy-and-hold жишгээр \$10{,}000 хөрөнгө +13.47\% буюу \$11{,}346.87 болсон.

Win rate 44.78\%--65.71\%, Expectancy ($R$) 0.79--1.63 хооронд эерэг гарсан (TP/SL=3:1 нөхцөлд break-even 25\%). Max drawdown 1.98\%--9.53\%, хамгийн урт дараалсан алдагдал 2--10. Execution rate 6.45\%--23.81\% мужид хэлбэлзсэн. Спред, гулсалтыг тусгасан боловч шимтгэл, свопыг бүрэн хамруулаагүй.

Инженерчлэлийн хувьд загвар сургалтаас эхлээд серверийн хэсэг (backend)-ээс харагдах хэсэг (frontend)-д хүргэх нэгдсэн урсгалыг прототипийн түвшинд хэрэгжүүлсэн.

\noindent\textbf{5. Ном зүй}

\small
Товчлолын гол эх сурвалжууд:
\begin{enumerate}
    \item Hevner, A. R., et al. (2004). Design Science in IS Research.
    \item Peffers, K., et al. (2007). A DSR Methodology.
    \item Chen, T., \& Guestrin, C. (2016). XGBoost.
    \item Ke, G., et al. (2017). LightGBM.
    \item Pardo, R. (2008). Trading Strategies.
    \item Bailey, D. H., et al. (2014). Backtest Overfitting.
\end{enumerate}
\normalsize

\end{multicols}

